%% Cover letter for the Nature Communications submission. Every number here also appears in the
%% manuscript; paper/check_numbers.py reads this file for that reason.
\documentclass[11pt]{article}
\usepackage[utf8]{inputenc}
\usepackage[T1]{fontenc}
\usepackage[a4paper, margin=2.5cm]{geometry}
\usepackage{hyperref}
\pagestyle{empty}
\setlength{\parskip}{0.6em}
\setlength{\parindent}{0pt}

\begin{document}

\hfill 12 September 2026

To the Editors, \emph{Nature Communications}

Dear Editors,

We submit ``Parameter-free fusion of histology, transcriptome and stage predicts bladder cancer
survival'' for consideration as an Article.

Multimodal models that combine whole-slide images with tumour transcriptomes are ranked for bladder
cancer almost entirely on one public benchmark, and every method paper on it that reports a clinical
baseline builds that baseline from tumour grade. In muscle-invasive bladder cancer 94.4\% of patients
share one grade level, so the covariate the field compares against cannot separate the patients it
does not distinguish. Rebuilding the same clinical arm with pathologic stage, from the benchmark's own
released files and through an identical pipeline, raises it from 0.5666 to 0.6638 concordance and more
than halves the apparent added value of three published multimodal constructions. The bias replicated
in an independent cohort where grade and stage both vary. This changes how a whole subfield's results
should be read, and it is a correction any group can adopt immediately.

We then ask what a method should look like at this sample size. With 359 patients and 113 events,
a ridge model fitted to pure noise reaches 0.7474 in sample at sixteen parameters, and a fusion
whose weights are fitted on the held-out fold is worth only $+0.008$ over equal weighting. ModRank
therefore fits one penalised Cox model per modality and fits nothing in the combination. It reaches
0.7212, the highest value reported on these folds, with 1,049 coefficients against 24.7 and 26.9
million for the two competitors whose code runs end to end. On identical inputs it exceeds a
concatenated Cox model and learned stacking in every one of 24 re-partitions of the same patients,
and it holds its concordance when whole hospitals are held out.

We also report what does not hold. Neither input-parity margin against a competitor reaches
significance at this cohort size, the protocol's own selection correction was mis-specified and the
result does not clear the corrected bar, and in three external cohorts without histology the
two-modality form of the rule does not exceed the clinical stage block beyond each cohort's noise
floor. Those boundaries are in the abstract and the Discussion rather than in a supplement, because
the paper's first finding is about how this field measures itself and it would be inconsistent to
report our own limits any less plainly.

The work is computational and uses public data only. Every number in the manuscript is regenerated
from committed result files by a checking script that is part of the archived code, and each analysis
ran under a protocol frozen before its outcomes were seen, including the external cohorts. There is
no preprint, and we have not opted in to any preprint-linked review service.

We declare no competing interests. The manuscript has not been submitted elsewhere.

Yours sincerely,

Guanxing Chen \\
City University of Hong Kong (Dongguan), Dongguan, China \\
\texttt{guanxing.chen@cityu-dg.edu.cn}

\end{document}
